% ================================================================
%  第 3 章 一封信 (前半: 原书 p.73-91, 边码 L1-L30)
%  源文: _work/src1/task-ch03a.md
% ================================================================
\chapter{一封信}

\noindent 1985年5月\par

% ---------------- p.73 ----------------
\section{(1) 千页之信}

\origpage{L1}我在这里寄给你的这份文本,是我所在大学张罗着打字、并只印了有限几册的,可它既不是抽印本,也不是预印本。它的名字——《收获与播种》——已经够清楚地预告了这一点。我把它当作一封长信寄给你——而且是一封再私人不过的信。我之所以寄给你,而不是满足于让你哪天(如果有这份好奇心)在书店里卖的某本书里看到它(倘若有哪个出版商疯到肯冒这个险……),是因为在它里面我更多是对你说话,而不是对别人。写它的时候,我不止一次想到你——得说,这封信我写了一年多,是把自己整个儿放进去写的。这是我给你的一份馈赠;写的时候我尽心竭力,把我(每一刻)能拿得出的最好的东西给出来。这份馈赠会不会被接受,我不知道——你的回答(或者你的不回答……)会让我知道……

在寄给你的同时,我也把《收获与播种》寄给数学界里所有那些同事、朋友或(旧日的)学生:我曾在某个时候与他们密切相连,或者他们以某种方式出现在我的反思里,点了名也好,没点名也好。你很有可能就在其中;而只要你用心去读,不只用眼睛和脑子去读,你必定会在那些没有点你名字的地方认出自己。我也把《收获与播种》寄给另外几位朋友,是科学家也好,不是也好。

% ---------------- p.74 ----------------
\origpage{L2}这封你正在读的“导引之信”,它向你预告并呈上一封“千页之信”(这还只是开头……),也将权作前言。写这几行的时候,那篇前言还没有写出来。此外,《收获与播种》由五个部分组成(还不算一篇“抽屉式”的导论)。我这里寄给你的是第一部分(《愚妄与更新》)、第二部分(《葬礼(1)——或中国皇帝的袍子》)和第四部分(《葬礼(3)——或四种运算》)\fnote{1}{}{我把那些以某种名义出现在我的反思里、但我并不亲自认识的同行另作别论。我只给他们寄去《四种运算》(那一部分与他们关系更近),同时寄去由这封信构成的“O 分册”,以及《收获与播种》的导论(外加前四部分全体的详细目录)。}。我觉得这几部分与你关系更近。第三部分(《葬礼(2)——或阴阳之钥》)无疑是我这份见证里最私人的部分,同时也是在我看来——比别的部分更甚——具有“普遍”价值的部分,超越它诞生时周遭的那些特殊境况。第四部分(《四种运算》)里,我这儿那儿提到这一部分;不过第四部分可以独立来读,甚至在很大程度上可以独立于前面三个部分来读\fnote{2}{(*)}{总的说来你会看到,每一个“节”(在《愚妄与更新》里)或每一则“按语”(在《收获与播种》其后三个部分的任何一个里),都有它自己的统一与自治。它可以脱离其余部分单独读,正如人看着一只手、一只脚、一根手指或一根脚趾、或整个身体上任何一大块一小块,也能看出兴味、得到乐趣,同时却不会忘记:那只是一个整体的一部分,而唯有这个整体(它始终留在未言说之中)才赋予它全部的意义。}。如果你读了这里寄去的东西,想给我回信(这也是我的愿望),或者它让你想读读缺的那一部分,那就告诉我。只要你回信让我感到你的兴趣超出了纯粹表面的好奇,我很乐意把它寄给你。

\section{(2) 《收获与播种》的诞生(一份闪电式的回顾)}

\origpage{L3}在这封前信里,我现在想用几页篇幅(如果办得到的话)告诉你,《收获与播种》里讲的是什么——说得比那个副标题本身更详细些:“一位数学家对往昔的反思与见证”(我的往昔,你猜得到的……)。《收获与播种》里东西很多,各色人等在其中看到的想必也很不同:一趟发现一段往昔的旅行;一场关于存在的冥想;一幅某个圈子、某个时代的风俗画(或者说,是从一个时代悄悄而不可抗拒地滑向另一个时代的图画……);一场调查(有时近乎侦探小说,另一些时候则近乎在数学这座大都市的底层上演的斗篷与长剑小说……);一次大规模的数学漫游(它会播下不止一颗……);一部应用精神分析的实用论著(或者,随你挑,一本“精神分析小说”);一篇关于自我认识的颂词;“我的忏悔”;一本私人日记;一种关于发现与创造的心理学;一份起诉书(照例是毫不留情的……),甚至是在“数学的美好世界”里算的一笔总账(而且不给人留半点情面……)。可以肯定的是,写它的时候我一刻也没有感到无聊,而与此同时我见识了、也看尽了各色各样的东西。要是你那些要紧的事务让你有空读它,那我倒奇怪你会读得无聊。除非是逼着你读,谁知道呢……

显然,它并不只是对数学家说话。不过也确有一些时候,它更多是对数学家说的,而不是对别人。在这封写给“《收获与播种》之信”的前信里,我却正想把那些可能与你、作为数学家,关系较近的东西概括出来、并着重凸显出来。
% ---------------- p.75 ----------------
\origpage{L4}要这么做,最自然的方式就是照直对你讲讲:我怎样一针引一线,接二连三写出了刚才提到的那四五块“砖头”。

如你所知,1970年我离开了数学的“大世界”,起因是我所属的机构(IHÉS)里一桩军费的事。此后几年,我搞反军备、搞生态,那种“文化大革命”式的行动,你大概在这儿那儿也听到过一些回声;然后我几乎从流通中消失了,天知道丢到了外省的哪所大学里。有传闻说,我把时间都花在放羊和打井上。实情是:除了许多别的营生,我还像所有人一样,硬着头皮去系里上课(这就是我那点毫无新意的糊口活儿,今天还是)。偶尔,甚至能连着几天、几星期、几个月,断断续续地重做一些数学——我手头满是装满我鬼画符的纸盒,能看懂的想必只有我自己。但那都是些与我从前做的事,至少乍看之下,极为不同的东西。1955年到1970年之间,我心爱的题目是上同调,尤其是各色簇的上同调(特别是代数簇)。我判断自己在这个方向上已经做得够多,尽可以让别人不用我也能应付;既然要做数学,那也该换一张唱片了……

1976年,一种新的激情进入我的生命,其强烈不亚于当年的数学激情,而且说来还是这激情的近亲。这就是对被我称作“冥想”的东西的激情(事物总得有个名字)。这个名字——正如任何别的名字在这里也免不了——必然会招来无数误解。与数学中一样,这是一件发现的工作。关于它,我在《收获与播种》里这儿那儿有所表述。总之,显然这里有的是活儿,够我干到死。而且,不止一次,我当真以为数学已成往事,今后我只会去管那些更严肃的事——以为我要去“冥想”了。

然而我最终还是不得不承认(那是四年前),数学的激情并未就此熄灭。甚至,连我自己也不太明白是怎么回事,并且暗自吃惊:我这个人(近十五年来)早已不再想在有生之年发表一行数学,却忽然发现自己被卷进了写一部数学著作的活儿里,那部书显然没完没了,还要一卷一卷地出;既然已经上了路,我索性要把自己以为在数学上有话要说的东西,统统倒进一套叫《数学反思》的(无穷?)系列书里,从此再没人提它。

那是两年前,1983年春。当时我已经太忙于写《追寻栈》(第一卷)了——那本该也构成《数学反思》(数学的)第一卷——顾不上追问自己身上发生了什么事。九个月后,理所当然地,这第一卷算是写完了,差不多也就是说,只剩导论要写、全稿要重读、注要加——然后就可以付印……

% ---------------- p.76 ----------------
\origpage{L5}直到此刻,那部书仍未完成——一年半以来它一动没动。那篇还欠着的导论,已经越过了(打字稿)一千二百页的大关;等它真正干完,少说也有一千四百页。你已经猜到了,这篇所谓“导论”,不是别的,正是《收获与播种》。据最新的消息,它该构成那套预定“系列”的第一、二卷,再加上第三卷的一部分。于是那套系列就此改名,干脆叫《反思》(不一定非得是数学的)。第三卷剩下的部分,主要由数学文本构成,如今这些对我比对《追寻栈》还要烫手。《追寻栈》就先等着吧,等明年再弄它的注、索引,当然还要一篇导论……

第一幕终!

\section{(3) 老板之死——弃置的工地}

\origpage{L6}我觉得,现在是时候做几点说明了:为什么我那么突然地离开了一个我在其中显然过得挺自在、过了二十多年的世界;为什么我起了那个古怪的念头,要“回来”(像个还魂的人……),而在我缺席的这十五年里,人家本来过得挺好;最后,为什么一部六七百页的数学著作的导论,竟写到一千二(或者一千四)百页。而在这里,一进入正题,我大概就要让你难受了(抱歉!),甚至惹你恼火。因为毫无疑问,你跟我从前一样,喜欢把你所属的那个圈子看得很“玫瑰色”——那个圈子里有你的位置、你的名头,诸如此类。我知道那是什么滋味……而这一回,恐怕要有点刺耳了……

我在《收获与播种》里这儿那儿谈到我离去的那段插曲,没有太在其中停留。那次“离去”在那里更像是我作为数学家的生命里一个重要的断口——我作为数学家的种种事件,都不断地相对于这个“点”来定位,作为“之前”与“之后”。要把我从一个我深深扎根的圈子、从一条早已深深划定的“轨迹”上拽开,需要一记极猛的冲击。这冲击来自碰撞:在一个我深深与之认同的圈子里,撞上了某种形式的腐败\fnote{3}{}{这里说的是各国科学家整体——以“建制”为首——与军事机器之间毫无保留的合作,把它当作经费、声望与权力的方便来源。这个问题在《收获与播种》里只顺带掠过一两次,例如去年4月2日那则《敬重》(第179号,第1221—1223页)。};对这腐败,我直到那时都选择闭上眼睛(仅仅做到自己不参与)。事后回看,我明白:在那事件之外,还有一股更深的力量在我身上起作用。那是一种强烈的、内在更新的需要。这样的更新,在一个地位显赫的科研机构那温吞的学术氛围里,是无法完成、也无法持续的。我身后是二十年紧张的数学创造、是二十年过度的数学投入——同时也是二十年漫长的、在“密闭罐”里的精神停滞……不知不觉间,我窒息了——我需要的是开阔处的空气!我那“天赐的”离去,标志着一场漫长停滞的骤然终结,也是我身上那些深刻力量走向平衡的第一步——那些力量曾被弯折、被拧紧在一种极度失衡、僵固的状态里……那次离去,真真正正是重新出发——是新旅程的第一步……

如我所说,我的数学激情并未就此熄灭。它在一阵阵零星的思考中找到了表达,走的路与我“从前”为自己划定的全然不同。至于我留在身后的那部作品,那部“从前”的作品——无论是白纸黑字印出来的,还是那部也许更为本质、却还没能找到书写之途、未能成书的作品——它大可以显得,而且在我眼里也确实显得,已经与我脱了钩。在去年之前(那是在《收获与播种》里),我从没起过念头,要在那些零零星星传回来的回声上稍作停留。我很清楚:我在数学里所做的一切,尤其在我1955到1970年那个“几何的”时期所做的,都是些本该被做出来的事;而我看见或瞥见的那些东西,都是本该出现、本该被拉到光天化日之下的东西。我也清楚:我做的工作,以及我让人做的工作,都是做得好的工作,是我整个人都投进去的工作。我把全部力气与全部爱都放了进去;而(在我看来)它从此是自治的了——一个活生生的、强健的东西——不再需要我来当它的母亲。在这方面,我离去时心里完全平静。我毫不怀疑:我留下的那些写出来的与没写出来的东西,我是把它们留在了可靠的手里;那些手会照看着它们展开、生长、繁衍,依照它们作为活物、作为强健之物的本性。

在这十五年紧张的数学工作中,在我身上绽放、成熟、长大的,是一种宏大的、统一的眼光,
% ---------------- p.77 ----------------
\origpage{L7}它化入少数几个极简单的有力观念。那眼光是一种“算术几何”的眼光,是拓扑学、几何学(代数的与解析的)与算术的综合,我在韦伊猜想里找到了它的第一个胚胎。正是它,成了那些年里我主要的灵感之源——在我看来,那些年尤其是这样一段时期:我从中提炼出这门新几何的主导观念,并打制出它几件主要的工具。这眼光、这些有力观念,对我来说已成了某种第二本性。(而在与它们断绝一切接触近十五年之后,我如今发现,这“第二本性”在我身上依然活着!)它们对我来说如此简单、如此自明,以至于理所当然地以为“所有人”都随时间与我一道把它们吸收、据为己有了。只是最近,就在这几个月里,我才明白:无论是那眼光,还是那几条一直是我恒常向导的“有力观念”,在任何一个已出版的文本里都找不到逐字的书写,至多不过是字里行间。而尤其要紧的是:我自以为已经传达出去的那眼光,以及承载它的那些有力观念,在达到完全成熟二十年后的今天,仍然不为任何人所知。我这个工匠,这些我有幸发现之物的仆役,同时也是它们在其中仍然活着的唯一一人。

\origpage{L8}我打制的这样或那样的工具,被这儿那儿用来“撬开”一个公认困难的问题,就像撬开一只保险箱。那工具显然结实。可我知道它还有另一种“力”,不只是万能钳子的力。它是一个整体的一部分,正如肢体是身体的一部分——那整体是它由之而来的,赋予它意义,它从那里取得力与生命。你尽可以用一根骨头(如果够粗)去砸开一个脑壳,这是明摆着的事。但那不是它真正的功能,不是它存在的理由。而我看见这些工具散落各处,被这个人那个人抓在手里,有点像被仔细剔净、剥光的骨头,是他们从一具身体上扯下来的——一具活的身体,而他们装作不认识它……

我在这里用字斟句酌的话说出的东西,是经过长久反思之后的结论;但它想必早已被我一点一点、含含糊糊地感知到了,年复一年,在那个尚未成形、还在寻求进入自觉的思想与形象、进入清晰分节之言说的未言说层面。我曾判定:这段过去,说到底与我无关了。那些远远近近传到我这里来的回声,尽管都被过滤过,却仍然意味深长,只消我稍作停留。我曾以为自己不过是众多工匠中的一个,在五六个正开足马力干着的“工地”上忙活\fnote{4}{}{关于这些被弃置的“工地”,我最后在三个月前那组《荒废的工地》(第176$'$至178号)里做了梳理。一年前,在发现“葬礼”之前,这事已经谈过一次,就在我重新与自己的作品接上联系、谈到它后来的遭遇的第一则按语《我的孤儿们》(第46号)里。}——也许是个更有经验的工匠,是那个长辈:当年他曾独自在这些地方干了许多年,直到一拨称心的接班人来了;是长辈,不错,但说到底与别人并无不同。可是你看,这一位一走,那就像一家石作行在老板突然去世之后宣告破产:几乎可以说一夜之间,工地就空了。“工匠们”走了,每人腋下夹走一点小零碎,觉得自家用得着。钱箱也没了,于是从此再没有任何理由让他继续费力干活了……

这同样是一种说法,是从一年多的反思与调查里澄清出来的。但毫无疑问,早在我离去后的最初几年,这已经是“在什么地方”被感知到的东西了。除了德利涅关于弗罗贝尼乌斯特征值绝对值的那些工作(就是那个“威望问题”,我最近才弄明白……)——每当我在远远近近碰上一位旧日的亲近同伴,曾与我同在工地上干过的,我问他“那么……呢?”,得到的总是同一个意味深长的手势:双臂往上一抬,像是在讨饶……显然,所有人都在忙着比那些曾放在我心头的事更重要的事;也显然,就在所有人摆出一副忙碌而要紧的神气忙活的时候,其实没做出多少东西。最要紧的
% ---------------- p.78 ----------------
\origpage{L9}那一块已经不见了——那种赋予局部任务以意义的统一,还有,我想,一种温度。剩下的是一堆从整体上脱落下来的零散任务,各人在自己的角落里孵着自己的小宝贝,或者好歹让它生点利息。

即便我想抗拒,眼看一切都戛然而止,我当然还是难受:再也听不到有人谈起 Motif,谈起 topos,谈起六种运算,谈起德拉姆系数,也听不到有人谈起霍奇系数,谈起那个“神秘函子”——它本该把各个素数对应的 $\ell$-进系数,围绕德拉姆系数,在同一把扇面上彼此联系起来;听不到有人谈起晶体(除非是得知它们始终原地不动),也听不到有人谈起“标准猜想”以及其他那些我提炼出来、显然代表着关键问题的东西。就连随《代数几何原理》(靠着迪厄多内不倦的协助)开始的那项宏大的基础工作——那几乎只消顺着已经取得的势头继续下去——也被撇在一边:所有人都满足于搬进别人耐心拼装、架设、擦亮了的墙壁与家具里。工匠一走,竟没有一个人想到轮到自己挽起袖子、动手拿起抹刀,去盖那些还等着盖的许多房子——那些好房子,能住人的,为自己,也为所有人……

我又一次忍不住,用完全自觉的形象把话接了下去,这些形象是从反思的工作里析出、浮上来的。但我毫不怀疑:这些形象早已以某种形式存在于我存在的深处。我一定早已感到,对我作品、同时也对我本人的那场“葬礼”,有着某种悄悄逼近的真实性;而这真实在去年4月19日突然强加于我,带着不可辩驳的力量,并且就叫那个名字:“葬礼”。相反,在自觉的层面上,我几乎不会想到要为此愤慨,甚至不会为此悲伤。说到底,不管是不是旧日的“亲近者”,那只是当事人自己的事:他选择拿自己的时间干什么。如果当年似乎推动过他、启发过他的东西不再启发他了,那是他的事,不是我的事。如果同样的事,以完美的一致,无一例外地落在我所有旧日的学生头上,那仍然是他们各自的事;我还有别的猫要抽,没工夫去追究这能有什么意思,就这样完事!至于我留下的那些东西,那些与我仍有某种深刻而不被察觉的联系的东西——即便它们显然被弃置在那片荒废的工地上,我自己却很清楚:它们不是那种会怕“时间的凌辱”、会怕风尚起伏的东西。如果它们还没有进入公共的遗产(当年我本以为已经进入了),它们迟早会在那里扎根,十年或一百年,说到底无关紧要……

\section{(4) 一阵葬礼的风}

可是,尽管这些年来我一直乐于回避那场大规模“葬礼”的含混感知,它却不依不饶地提醒我记起它,以别的面貌,而且不那么无害——不只是对一部作品的冷淡而已。我渐渐得知——我也说不清是怎么得知的——那被遗忘的眼光里的好几个概念,不仅已经废弃不用,而且在一个体面的圈子里,成了居高临下鄙夷的对象。情形尤其如此的是那个关键的统一性概念:topos(广义空间)——它正处在新几何的核心,正是它提供了拓扑学、代数几何与算术共有的几何直觉,也正是它让我既能提炼出平展的、$\ell$-进的层论工具,又能提炼出晶体上同调的那些主导观念(诚然,此后它们多多少少被遗忘了……)。说实话,
% ---------------- p.79 ----------------
\origpage{L10}连我自己的名字,也年复一年地、悄悄地、莫名其妙地成了取笑的对象——成了没完没了、稀里糊涂的“轰炸”的同义词(正是针对那些著名的 topos,或者那些“Motif”——他拿这些玩意儿把你耳朵都磨出茧子来了,可谁也没见过……),成了长达千页地把一根头发劈成四瓣的同义词,成了关于那些人人都一向知道、根本不用等他说的事滔滔不绝、臃肿巨大的唠叨的同义词……大致就是这类调子,不过是压低了声音,靠言外之意,带着“在高蹈之士与高雅之交之间”应有的那份体贴。

在《收获与播种》里进行的那场反思中,我想我已摸到了在各人身上起作用的那股深刻力量——它就藏在那副取笑与俯就的神气背后,面对一部其分量、其生命、其气息都非他们所能企及的作品。我还发现(撇开我这个人身上那些给我的作品与命运打下印记的特殊之处)那个秘密的“催化剂”,它促使这些力量以这种形式显现出来:面对一个完好无损的创造力所给出的意味深长的迹象,露出那种漫不经心的轻蔑。总之,在这场被取笑和轻蔑捂得悄无声息的葬礼上,他就是那位大主祭。说来奇怪,他同时也是所有人当中离我最近的那一个——而且是唯一一个曾经领会过某种眼光、并把它据为己有的人,那眼光里满是生命与强烈的力量。不过我又抢先了……

说实话,那些这儿那儿传回我这里的“隐约取笑的小阵风”,并没有特别伤到我。直到三四年前,它们在某种意义上还是无名的。我当然从中看到一个不大令人愉快的时代征兆,但它们并没有真正质疑我,也没有在我心里引起焦虑或不安。相反,更直接触动我的,是那些与我本人拉开距离的迹象,它们这儿那儿来自我昔日数学界里的不少朋友——尽管我离开了那个曾与我们共有的世界,我仍觉得自己与他们由同情的纽带相连——除了共同的激情与一段共同的过去所造就的那些纽带之外。在这方面,每一次我固然都难受,可也并没有太往心里去,而且(就我记忆所及)从没想过要把这三串迹象联系起来:被弃置的工地(以及被遗忘的眼光)、“取笑之风”,以及不少旧日朋友拉开距离。我给他们每个人都写过信,没有一个人回信。再说,从那以后,我写给旧友或学生、谈那些搁在我心头的事的信,常常没有回音,这已经不稀奇了。新时代,新风尚——我又能怎样?我只好不再给他们写信。可是(如果你是其中一位)我现在正在写的这封信,将是那个例外——一番再次向你递上的话——这一次接不接受,你自己看吧……

\origpage{L11}某些旧友与我本人拉开距离的最初迹象,如果我没有弄错,要追溯到1976年。也正是那一年,另一“串”迹象开始出现,那是我在回到《收获与播种》之前还要谈的。说得更准确些,这后两串迹象是同时出现的。就在我写下这些的此刻,我觉得它们其实不可分割,说到底,它们是同一现实的两个不同侧面或“面孔”,那一年闯进了我亲身经历的场域。至于我方才正要谈的那个侧面,那是一种系统性的“拒不受理”,\dmark 不声不响、无可辩驳,由一个“毫无裂缝的共识”\fnote{5}{}{这个“毫无裂缝的共识”,在《愚妄与更新》里这儿那儿偶尔被提到,到下一部分《葬礼(1)》里才终于成为一份详尽的见证与一番反思的对象,其中有“第X支队列”即“灵车”,由那些“棺材按语”(第93—96号)以及按语《掘墓人——或整个教团》组成。后者为《收获与播种》的这一部分收尾,同时也构成反思这“第二次吐息”的第一个成果。}专门留给1970年以后那少数几个学生及其同类——他们凭着各自的工作、工作方式与灵感,清清楚楚带着我的影响的印记。
% ---------------- p.80 ----------------
\origpage{L12}也许正是借这个机缘,我第一次察觉到那阵“隐约取笑的气息”,它通过他们,针对的是某种数学风格与某种数学进路——一种(按照当时在数学建制里显然已成普遍的共识)本不该存在的风格与眼光。

这里同样,是某种在无意识层面上被清楚感知到的东西。就在同一年,同一个荒诞的剧本(演示的是:一篇显然出色的博士论文竟无法发表)连续重复了五次,带着马戏团噱头般执拗的滑稽,它最终还是硬闯进了我自觉的注意。如今回想,我明白:当时某种现实正以一种善意的执拗在“向我示意”,而我在装聋:“喂,睁眼看看吧,大个子,留神点你眼皮底下正发生的事,这事跟你有关,真的……!!”我稍稍振作了一下,看了一眼(就那么一小会儿),半是发懵,半是心不在焉:“哦,是啊,好吧,有点怪,看来这儿是有人跟谁过不去,有什么东西横竖过不去,而且还这么齐整,说实话简直难以置信!”

这事简直难以置信到了这个地步,以至于我赶紧把那个噱头、把那场马戏都忘了。确实,我并不缺别的有趣的事可做。可这并没挡住马戏在随后几年里继续提醒我记起它——不再是噱头的调子,而是暗中以羞辱为乐的调子,或者照准脸上抡一拳的调子;只不过这里的人物都是些高雅人士,那一拳在这里也就必然采取更高雅的形式,但一样有效,全凭那些高雅人士的创意……

我感受到如同“照准脸上抡来的一拳”(是别人挨的)的那个事件,发生在1981年10月\fnote{6}{}{这段插曲见按语《棺材之三——或有点过于相对的雅可比簇》(第95号),尤其第404—406页。}。那一次,而且是从种种新精神的迹象不断向我袭来以来的第一次,我被打中了——大概比拳头落在我自己身上还要更重;那一拳是别人挨的,而那个人是我所疼爱的。他多少有点像个学生,而且是一位天分非凡的数学家,刚做出了一些漂亮的东西——不过这是个细节,说到底。不是细节的,是当时我有三个“从前的”学生,直接与那件事站在一边,而当事人(并非没有道理地)把那件事当作一种羞辱与侮辱。另外两个旧日的学生,早已有机会居高临下地对待过他,像有钱人打发一个穷酸书生那样\fnote{7}{}{在上一则脚注所引的那则按语里顺带谈到过。}。还有一个学生,三年后也会跟上同样的步子(而且依旧是“照脸一拳”的风格)——不过这一点我当时当然还不知道。当时叫住我的东西,已经绰绰有余。仿佛我那段从未被审视过的数学家的过去,忽然借着我五个旧日学生中某些人的嘴脸,朝我做了一个狰狞的鬼脸——这些人如今都成了要人,有权有势,而且倨傲……

那本该是再也合适不过的时机,去停下来,去探问那突然以如此暴力向我招呼的东西究竟是什么意思。可在我身上的某个地方早已决定(而这事从来不必说出口……):那段“从前”的过去,说到底与我无关了,我没有必要在那里停留;如果它现在似乎用一种我认得再清楚不过的声音——轻蔑时代的声音——在招呼我,那一定是认错了人。然而,我连着几天、也许几周,被焦虑捆得紧紧的,却连承认它都不肯。(只是去年,借着写《收获与播种》使我重回到这段插曲,我才终于得知有这份焦虑,它一出现就被控制住了。)我没有去确认它、探问它的意义,而是四处活动,我写信给
% ---------------- p.81 ----------------
\origpage{L13}左右各方,写“那些非写不可的信”。当事人甚至还费心给我回了信,自然是些含糊其辞的信,根本不触及任何实质。浪头终于平息,一切恢复秩序。直到去年,我几乎没再想过这事。可这一次,它留下了一道伤,或者更像一根扎人的刺,让人避着不去碰;一根养着那道伤口的刺,而那伤口只等着愈合……

这大概是我作为数学家的生涯里经历过的最痛苦、最难熬的经验——当我有幸看到(却又不肯真正去认清我眼睛所看见的):“某个我疼爱过的旧日学生或同伴,以悄悄压碎另一个我疼爱的人为乐,并且在那人身上认出了我”。它当时给我的打击,肯定比去年我那些相当疯狂的发见还要重——那些发见(在浮浅的目光看来)也许显得更加不可思议……确实,这经验在同样的调子上引起了另外几次共鸣,只是没那么暴烈,当时也就有点“漏过去了”。

\origpage{L14}这也让我想起:同样是1981年,也是我与旧日学生中唯一一个在我离去后仍保持经常联系的人的关系发生急剧转折的一年;而这个人,十五年来在数学层面上一直是我“优先的对话者”。正是那一年,已经出现几年的“那些表示轻蔑的迹象”\fnote{8}{}{这段插曲见按语《两个转折》(第66号)。}“突然变得如此粗暴”,于是我中断了与他的一切数学交流。那是在前面那件“照脸一拳”的事之前几个月。事后回看,这巧合让我觉得触目;但我不认为当时自己做过任何联系。那些事被归进一个个“分开的抽屉”;何况有人已经宣布过,这些抽屉其实无关紧要——案子已经了结!

这也让我想起:还是这1981年6月,已经开过一场颇为精彩的研讨会,在不止一个意义上值得记住——那场研讨会着实配得上以“反常讨论会”这个抹不掉的名字进入历史(或者说进入历史剩下的那点东西……)。我是在去年5月2日认识它的(更确切地说,是它朝我砸下来的!),就在发现那有血有肉的“葬礼”(4月19日)两周之后——我立刻明白,我撞上了“封神”。当然,是一场葬礼的封神,但同时也是一种轻蔑的封神——轻蔑的对象,是我们这门科学存在两千多年来数学家伦理那条心照不宣、不可动摇的根基:那条最基本的规矩,即不得把从别人那里拿来的想法与结果当作自己的东西拿出来。而我此刻注意到这两个事件在时间上那引人注目的巧合——它们看上去性质与分量都大不相同——不禁愕然:这里显露出的,正是对一个人的敬重与对一门技艺、一门科学的基本伦理规则的敬重之间那种深刻而明显的联系;正是这些规则,使这门技艺的从事不成为一场“混战”,使那些以擅长此道著称、并在此定调的人整体上不成为一个毫无顾忌的“黑手党”。不过我又抢先了……

\section{(5) 旅行}

我想,关于我的“重返数学”、并一针引一线地关于《收获与播种》的写作所处的来龙去脉,我大致已经绕了一圈。是去年3月底,在《愚妄与更新》的最后一节(《一段过去的重量》,第50号)里,我才终于想到要追问这次意外重返的缘由与意义。说到“缘由”,最强的一条想必是那种既含混又迫人的感觉:那些强健有力的东西,我曾
% ---------------- p.82 ----------------
\origpage{L15}当年托付给一些充满爱意的手——“它们却在坟墓里,与风、雨、阳光的恩泽隔绝,在我看不见它们的这十五年里萎坏着”\fnote{9}{}{引自按语《坟前的旋律——或自得》(第167号),第826页。}。我一定是渐渐地明白了,而且此前从未想到要对自己说出来:最终撬开那些虫蛀的木板、放出那些被囚禁的活物的,不是别人,正是我自己;那些活物生来不是为了让它们在封闭的棺材里烂掉,而是为了在开阔的空气中绽放。而那些围绕这些包了衬垫、臃肿不堪的棺材(无疑是照着那位可敬的死者的模样……)的假作哀戚与暗含取笑的神气,想必也“终于在我心里唤醒了一根近十年来有点睡过去的斗志之弦”,以及投身混战的欲望……\fnote{10}{}{见《一段过去的重量》(第50节),尤其第137页。(**)。}。

就这样,两年前,一件本来只打算作几天、顶多几周的快速踏勘——对某一处被撇下的“工地”的探看——变成了一部N卷本的数学长篇连载,插进了那套著名的新丛书《反思》(“数学的”,等哪天把那个多余的限定词删掉再说)。而且,从我得知自己正在写的是一部要出版的数学著作的那一刻起,我也就知道:除了或多或少合乎惯例的“数学”导论之外,我还要给它加另一篇“导论”,性质更私人。我觉得重要的是,我要解释我的“重返”——那根本不是重返一个圈子,而仅仅是“重返”紧张的数学投入,重返发表我笔下写出的数学文本,为期不定。同样,我也想解释清楚我现在写数学时所怀的精神,它在某些方面与我离去前的写作精神大为不同——那是一种发现之旅的“航海日志”精神。更何况,还有些别的事搁在我心上,它们无疑与这些相关,但我觉得更为本质。我早已认定,我要慢慢来,把该说的话说出来。那些事还含混不清,对我来说,却与我将要写的这些卷书、以及它们将插入其中的《反思》的意义不可分割。绝不该把它们草草塞进去,像为占用了匆忙读者宝贵的时间而道歉似的。如果说《追寻栈》里有些东西值得他、也值得所有人知道,那恰恰是我留待在这篇导论里说的。要是二三十页不够把它们说完,我就写上四十页,甚至五十页,那又何妨——何况我并没有强迫谁来读我……

\origpage{L16}《收获与播种》就是这样诞生的。我1983年6月写下了预定那篇导论的最初几页,那是在写《追寻栈》第一卷中间的一段空档。然后去年2月我又接着干,那时我那卷书已经基本完成好几个月了\fnote{11}{}{其间我花了足足一个月,思考一个伪直线系统的“结构曲面”,那是用一条伪直线相对于这样一个系统的所有可能“相对位置”的全体来表示的。我还写了《纲领概要》,它将收入《反思》第三卷。}。我本以为这篇导论会是个机会,让我把心里还有两三处略微模糊的东西弄清楚。可我丝毫没料到,它会像我刚写完的那卷书一样,成为一趟发现之旅;而且这趟旅行所进入的世界,比我准备在那卷已成的书和随后几卷里踏勘的世界要丰富得多、广阔得多。日子一天天、一周周、一月月过去,我也不太清楚发生了什么,这趟新的旅行就这样继续下去:去发现一段过去(三十多年来被执意回避的……),发现我自己,发现我与这段过去的种种联系;也去发现数学界里某些曾与我亲近、而我了解得如此之差的人;最后甚至,在
% ---------------- p.83 ----------------
\origpage{L17}这股势头之中,还成了一趟数学发现的旅行——因为十五二十年来,我第一次有闲情回到某些我离去时留得滚烫的问题上去\fnote{12}{}{五六十年代,我常常压抑自己想投身这类油水十足、滚烫的问题的欲望,因为被没完没了的基础工作占住了;那些工作,没有人能、也没有人愿意替我接着做,而我离去之后,也没有人有心把它继续下去……}。总之我可以说,我在《收获与播种》的篇幅里进行的,是三趟彼此紧紧交织的发现之旅。而这三趟,没有一趟随着第一千二百来页那个句号而告终。我这份见证将要收拢的回声(甚至包括那由沉默构成的回声……)都将成为这趟旅行的“续篇”。至于它的终点,这趟旅行肯定属于那种从未被走到头的旅行——也许,甚至到我们死的那天也走不到头……

于是我终于又回到了出发点:如果办得到,事先告诉你,《收获与播种》里“讲的是什么”。可也真是,前头那些篇幅,即使我没去找,也已经多多少少告诉你了。也许更有意思的做法是顺着势头往下讲,而不是去“预告”。

\noindent 1985年6月\par

\section{(6) 阴影的一面——或创造与轻蔑}

前面那些篇幅,是上个月借着一段短促的“空档”写下的。其间,我终于把《四种运算》(《收获与播种》的第四部分)最后收拾停当——现在只差把眼前这封信、这封“前信”也写完(它同样摆出一副要长到无法无天的架势……),一切就终于可以付排、可以复印了。我都不再相信这事了,快一年半了,我一直“就要写完”那批出了名的按语!去年2月(还有前一年的6月),当我动手写这篇对一部数学著作来说性质有点不寻常的“导论”时,(我想)主要有三类东西是我当时想说的。首先,我想说明我重返数学活动的意图,说明我写《追寻栈》这第一卷(我刚宣布它写完)时所怀的精神,也说明我打算带着《反思》去进行一趟更广阔的数学踏勘与发现之旅时所怀的精神。从此对我来说,不再是为某个正在孕育中的数学新世界呈上一套细致入微、一丝不苟的基础。那更该是些“航海日志”,工作在其中日复一日地继续,什么也不隐瞒,就跟它真实进行的样子一样,带着它的卡壳与失手,带着它执拗的回头,也带着它突然的跃进——一件日复一日被不可抗拒地拉着往前的工作(尽管有无数意外与不测),仿佛被一根看不见的线牵着——被某种飘忽、执着而可靠的眼光牵着。一件常常在摸索中前行的工作,尤其在那些“敏感的时刻”:某种还没有名字、还没有面孔的直觉,几乎无法察觉地浮上来;或者在某个新旅程的起点,应着最初几个想法与直觉的召唤、追着它们走——它们往往飘忽,不肯被语言的网眼捉住,而恰好是那种能够轻柔地捉住它们的语言,常常还欠着。首先要做的,正是让这样一种语言从一片显然空无的、摸不着的雾霭中凝结出来。那还只是被预感、尚未被瞥见、更谈不上被“看见”和被手指触到的东西,一点点地从不可秤量之物中澄清出来,从它的阴影与雾霭的外套里脱出来,取得形状、血肉与重量……
% ---------------- p.84 ----------------
\origpage{L18}工作的这一部分,样子不起眼,甚至(许多时候)可以说干得糟糕,却同时也是最微妙、最本质的部分——真正说来,正是在那里,某种新东西出现了,靠着一种紧张的注意、一种关切、一种对这脆弱得无限微妙、即将诞生之物的敬重。这是创造的部分,是受孕与在养育之母体温暖的幽暗中缓慢孕育的部分:从那看不见的双重始基细胞起,变成无形的胚胎,又日复一日、月复一月地,在一种幽暗而紧张、看不见也无形状的劳动中,变成一个血肉俱全的新生命。

这也正是发现工作中“幽暗的”部分,那“阴”的或“女性”的部分。与之互补的一面,那“明亮”的、“阳”的或“男性”的部分,则更像抡起锤子或大槌干活,干在磨得锋利的凿子上,或干在钢火淬好的楔子上。(那是些现成可用的工具,效力早已得到验证……)两个方面各有其存在的理由与功能,彼此不可分割地共生——或者说得更好些,它们正是那对不可分离的原始宇宙力量夫妇中的妻与夫;它们不断更新的拥抱,不断重新唤起那幽暗的创造劳作:受孕、孕育与诞生——孩子、新生之物的诞生。

\origpage{L19}在那篇出了名的、私人的、“哲学的”导论里,我觉出需要说的第二件事,恰恰是关于创造工作的本性。多年来我已明白:这本性一般都被忽视,被随手拈来的套话、被祖祖辈辈的压抑与恐惧掩盖了。它究竟在多大程度上确实如此,我此后才渐渐发现,日复一日、月复一月,在《收获与播种》里进行的那场反思与“调查”的全过程中。正是在这场反思的“发令”之初,在1983年6月写下的那几页里,我第一次被一个外表平淡、实则令人目瞪口呆的事实攫住——只消稍微停留一下就会看到:我刚才说到的发现工作中那“最富创造性”的部分,在那些本该呈现这样一种工作(至少呈现它最摸得着的果实)的文本或言谈里,几乎哪儿都不露面;无论是教科书和其他教学文本,还是原创论文与论著,还是口头课程与讨论班报告等等。看来几千年来,自从数学与其他技艺、科学诞生之初起,就存在着某种“沉默的合谋”,围绕着这些“不可告人的劳作”——它们是一切新观念(或大或小)绽放的前奏,那些观念更新着我们对这个世界某一部分的认识,而我们就在这个不断创造着的世界里活着。

说白了,看来对那个侧面、那个阶段——一切发现工作(以及一般创造工作)中最关键的阶段——的认识所受到的压抑,竟如此有效,如此被那些明明亲身了解这种工作的人内化,以至于人们常常会发誓说:即便这些人,也已经把关于它的一切自觉记忆的痕迹都铲除了。有点像在一个清教到极点的社会里,一个女人会从自己的记忆里铲除掉:与每一个她自认为有责任去擤鼻涕、擦屁股的孩子相关的,那使她受孕的一刻(她勉强承受的拥抱),那漫长的几个月妊娠(被当作一件不成体统的事),以及那漫长的几个钟点分娩(被当作一场有点恶心的苦难,最后总算解脱)。

这个比喻也许显得夸张,而且,如果我把它用到我今天还记得的、我二十年前自己所属的那个数学圈子的精神上,也许确实夸张。可在我于《收获与播种》里的反思过程中,我能够认识到——尤其在这最后几个月(随着《四种运算》的写作)以触目的方式认识到——自我离开数学舞台以来,在今天那些我曾熟悉的圈子里、并且(在我看来,至少在很大程度上)在整个数学界里成为律令的那种精神,发生了令人目瞪口呆的败坏\fnote{13}{}{何况这种败坏绝不限于“数学世界”。在整个科学生活里都能看到它,而在科学生活之外,在当代世界的全球尺度上也能看到。朝这个方向所作的初步确认与反思,见开启阴阳之思的那则按语《肌肉与内脏》(第106号)。}。甚至可能,无论凭我极为特殊的数学个性,还是凭围绕我离去时的那些境况,我那次离去都在一场本已在进行的演变中起了催化剂的作用\fnote{14}{}{这就是上页脚注所引那则按语里考察的演变。它与(我本人和我作品的)“葬礼”之间的联系,在《阴的葬仪(阳葬阴之四)》《天赐的时机——或封神》《不认(一)——或召回》《不认(二)——或变形》(第124、151、152、153号)诸按语里出现并得到考察。另见较晚的(在 RS IV 里的)按语《无用的细节》(第171(v)号,(c)部分“那些什么也不像的东西——或干枯”)与《家庭相册》(第173号,c 部分“那唯一者——或默许”)。}——那场演变,我当时丝毫未察觉(我那些同事与朋友里也没有人察觉,也许只有克劳德·谢瓦莱(Claude Chevalley)是唯一的例外)。我在这里尤其想到的这种败坏的那个侧面(它只是众多侧面中的一个\fnote{15}{}{在《收获与播种》里最常处在注意中心的侧面,尤其是在两个“调查”部分(RS II,即“中国皇帝的袍子”;RS IV,即“四种运算”)里,而且也许正是最让我“目瞪口呆”的那个侧面,是这一行当伦理的败坏:它表现为肆无忌惮的掠夺、贬损与舞弊,行之于当下某些最负盛名、最出色的数学家之间,而且(在很大程度上)是在众目睽睽之下。至于另外一些更微妙的侧面——它们其实与这一侧面直接相关——我请读者去看已经引过的那则按语(第173号 c 部分)《那些什么也不像的东西——或干枯》。})是对以下这一切心照不宣的轻蔑,乃至毫不含糊的取笑:在数学里,凡是不像锤子砸在铁砧上或凿子上那种纯粹劳作的——那些最微妙的创造过程(而且往往样子最不起眼)所受的轻蔑;一切灵感、梦、眼光(无论它们多么有力、多么丰饶)所受的轻蔑;乃至(推到极处)一切观念所受的轻蔑,无论它被构想、被表述得多么清楚:凡是没有白纸黑字写出来、没有以纯粹而硬邦邦的命题形式被写出、被登记、被编目、可以进入那些灌满我们巨型计算机无穷记忆的“数据库”的东西。
% ---------------- p.85 ----------------
\origpage{L20}(借用 C. L. 西格尔(C. L. Siegel)的一个说法\fnote{16}{}{这个说法在上一页脚注所引的那则按语里被引用并加以评论。})曾经有过一次非同寻常的“压平”,数学思想的一次“收缩”:它被剥去一个本质的维度,剥去它全部的“阴影的一面”、那“女性”的一面。诚然,由于一种祖传的传统,发现工作的这一面在很大程度上一直隐而不显,可以说从没有人谈过它——但与梦的深处源头(那些滋养伟大眼光与伟大构想的源头)的活的接触,此前(据我所知)从未丧失过。看来我们从今起已经进入了一个干枯的时代:在那里,这一源头不是被从中心处切断,而是通往它的路被堵死了,凭的是普遍轻蔑那不容上诉的判决,以及嘲弄的那些代表们。

我们似乎正在接近这样一个时刻:在每个人身上,不仅将被铲除一切与源头相近的工作、那“女性气质”的工作(被讥为“稀里糊涂”、“软塌塌”、“不牢靠”——或者在另一极端,被讥为“琐屑”、“小儿科”、“瞎轰炸”……)的记忆,而且这工作本身连同它的果实也将被连根拔除:那构想、酝酿并生出新概念与新眼光的所在。那也将是我们的技艺之实践被缩减为干巴巴、空洞无物的脑力“举重与哑铃”表演的时代,缩减为在竞赛中“攻破”(“难度有口皆碑”的)问题的种种壮举不断加码的时代——那是一个“超猛男”式的、亢奋而贫瘠的肥大的时代,接在三百多年创造性的更新之后。

\section{(7) 敬重与坚毅}

可是我又离题了,预先说出了反思教给我的东西。我原本是从双重的意图出发的,这双重的意图在反思开始之前就已清清楚楚在我心里:一是作一份“意图声
% ---------------- p.86 ----------------
\origpage{L21}明”,二是(与前者紧密相连,正如刚才显出的那样)就创造工作的本性说点什么。然而还有第三个意图,它在自觉层面上想必没那么清楚,却回应着一种更深、更本质的需要。它是由那些有时令人不知所措的“招呼”引起的,那些招呼从我的数学家的过去向我传来,借着那些曾是我的学生或朋友的人的声音(或者至少是其中不少人的声音)。在表皮层面上,这需要表现为一种想“把肚子里的东西倒空”、想说几句“不中听的实话”的欲望。但更深处,肯定是一种终于要去认识一段过去的需要,而那段过去我此前一直选择回避。首先正是从这种需要里,生出了《收获与播种》。这场漫长的反思,就是我对自身这种求知冲动、对那从外部世界、从我头也不回地离开的“数学世界”不断重新传来的招呼,日复一日作出的“回答”。除了《愚妄与更新》最开头那几页——构成其最初两章(《劳动与发现》与《梦与做梦者》)的那几页——之外,而且从接下《恐惧的诞生》(第18页)的那一章起,伴着一份根本不在计划之内的“见证”,我认为,正是这种认识我的过去并完全承担它的需要,成了写作《收获与播种》时起作用的主要力量。

\origpage{L22}从数学家的世界传到我这里、并在《收获与播种》全过程中(尤其是在第二、第四部分所进行的“调查”中)以新的力量反复落回我身上的那种招呼,一上来就戴上了自得的面具,要不就是(经过“精细调配”的)鄙夷、取笑或轻蔑的面具:有时是针对我,而(更主要地)是针对那些竟敢受我启发的人(他们当然没有料到等着他们的是什么),那些被某道心照不宣、不可更改的判决“归了类”,算作与我有一手的人。于是我又一次看见那条“明显”而深刻的联系在此显现:对他人这个人的敬重(或缺乏敬重);对创造之举、对其某些最微妙最本质的果实的敬重;最后还有对科学伦理那些最明显的规则的敬重——那些规则扎根于对自己和他人最基本的敬重,我忍不住要称之为我们这一行当里的“体面规则”。这些无疑都是某种基本的、本质的“自我敬重”的各个侧面。如果我试着用一句斩钉截铁的话,来总结《收获与播种》就某个曾是我的世界、我曾与之认同二十多年的圈子所教给我的一切,我会说:那是一个失去了敬重的世界\fnote{17}{}{这里同样,这个说法不只适用于某个有限的圈子——在那里我有充分机会就近看到这事——我觉得它概括了当今整个世界某种普遍的败坏。(可与第 L19 页脚注比较。)在《收获与播种》里所进行的“调查”之总结这一更有限的框架内,这个说法出现在去年4月2日那则《敬重》(第179号)里。}。

这件事在此前那些年里已经强烈地被感觉到,即便没有被说出来。它只是不断得到确认、变得明确,总是出人意料,有时令人目瞪口呆,贯穿《收获与播种》的全过程。它清楚显现在这样一个时刻:一种“哲学的”、一般的反思,忽然变成了个人的见证(就在开启前面提到的《恐惧的诞生》那一章的《受欢迎的异乡人》一节(第9号,第18页)里)。

然而这种感知并不是以尖锐或苦涩的怨责的调子出现的,而是(凭写作的内在逻辑、凭写作所唤起的另一种态度)以一种追问的调子出现的:在这场败坏里,在今天我所确认的这种敬重的丧失里,我自己那一份是什么?正是这个追问,贯穿并支撑着《收获与播种》这第一部分,直到它最终化解为一个清楚而毫不含糊的确认\fnote{18}{}{见《竞技数学》与《旋转木马收摊了》两节(第40、41号)。}。在此以前,这场败坏对我来
% ---------------- p.87 ----------------
\origpage{L23}说,像是忽然“从天而降”,无法解释,因而更加令人愤慨、难以容忍。而在反思的过程中,我发现它一直悄悄地进行着,想必没有人能在自己周围或自己身上察觉它,贯穿整个五十与六十年代,也包括我本人。

对这个卑微事实的确认——它想必明显得很,平淡无奇——标志着见证中第一个关键转折,以及一种立竿见影的质变\fnote{19}{}{次日,见证便深化为一场对我的沉思,并在随后的几周里保持这种特殊的品质,直到《收获与播种》这“第一次吐息”结束(以《一段过去的重量》一节,第50号)。}。这是我关于自己数学家的过去、关于我自己所要学的第一件本质的事。这种对我在普遍败坏中所负责任的认知(其敏锐程度随反思的时刻而或深或浅),贯穿《收获与播种》的全过程,成了一条底色、一个提醒。尤其当我的反思显出是在调查一个时代的失宠与不公时,情形更是如此。与理解的欲望、与那种驱动并推动一切真正的发现工作的好奇心一道,正是这种卑微的认知(它一再在途中被忘掉,又不顾一切地浮上来,而且总在最料不到的地方……)使我的见证免于(我相信)沦为对世界忘恩负义的干瘪怨责,甚至免于沦为与某些曾是我的学生或朋友(或两者兼是)的人“算总账”。这种对自己的不徇情,也给了我那种内心的平静,或者说那种坚毅,使我免于徇情于人的陷阱,哪怕只是免于一种虚假“谨慎”的陷阱。凡是我认为在反思的某一刻该说的话——无论关于我自己,还是关于我某位同事、某个旧日学生或朋友,还是关于一个圈子、一个时代——我都说了,从不需要去硬压自己的迟疑。至于那些迟疑,每次我只消仔细审视它们,它们便消失得无影无踪。

\section{(8) “我的亲近者”——或默契}

\origpage{L24}这封信里,我不打算把《收获与播种》写作中(或它某一阶段中)的所有“要紧时刻”(或所有“敏感时刻”)逐一巡览\fnote{20}{}{关于《收获与播种》前三部分全体的一个简短回顾与总结,你会在这两组按语里找到:《暮熟的果实》(第179—182号)与《发现一段过去》(第183—186号)。}。我只想说:这件工作里有四个界限分明的大阶段,或者说四口“气”——像呼吸的几口气,又像一道涌起的浪涛里接连而来的几道浪;那道浪涛我也不知是怎么从一片广袤、无声、静止而又流动、无界无名的大海里涌出来的,那海便是“我”;或者更确切地说,是一片无限辽阔、无限深邃的大海,远过于它所承载、所养育的那个“我”。这几口“气”、这几道“浪”,化成了《收获与播种》目前已经写成的四个部分。每一道浪来了,都不是我呼唤来的,也丝毫未曾预料;而任何时候我都说不清它要把我带到哪里,又何时结束。而当它结束了、新的一道浪已经接了上来,我还总以为自己在某一段冲势的收尾(那也终将,在最终的最终,是《收获与播种》的收尾!),其实我早已被同一个宏大运动的另一口气抬起来、托着往前走了。只有事后回看,这运动才清楚地显出,才毫不含糊地显出:在那被当作行动、被当作涌动而经历的东西里,有一个结构。

而且,那运动当然没有随着我在《收获与播种》上打的那个句号(完全是暂时的!)而结束,
% ---------------- p.88 ----------------
\origpage{L25}也不会随着我写给你的这封信的句号而结束;这封信也是那运动的一个“时段”。它也不是在某一个1983年6月或1984年2月的某一天诞生的——那天我坐到打字机前,要写(或接着写)某部数学著作的某篇导论。它诞生了(或者更确切地说,它重又诞生了……),就在冥想进入我生命的那一天……

可是我又离题了,任由即时的形象与联想把我带走(甚至卷走……),而不是乖乖守着一条“话头”、一条预定的线索。我今天本来是要接着讲,尽可能简短地讲去年4月那场“发现葬礼”的故事:当时我已有两个星期相信《收获与播种》已经写完——而就在那三四个星期的工夫里,种种发现怎样一个接一个、一个比一个更大更不可思议地朝我劈头砸下来——大得、疯得,以至于此后好几个月,我都极难“相信我健全官能的证词”,极难摆脱面对明证时那种暗暗的难以置信\fnote{21}{}{我试着用那个故事《中国皇帝的袍子》来表达这种困难,见同名按语(第77号);后来又在那则《尽到的义务——或真相的时刻》(第163号)里回到这一点。}。这种隐秘而顽固的不信,直到去年10月(在发现“光彩夺目的葬礼”六个月之后)才终于消散,那是在我的朋友、旧日学生(不错,是隐形的)皮埃尔·德利涅(Pierre Deligne)来我家拜访之后\fnote{22}{}{这次拜访的叙述,见我刚引的那则按语(在上页脚注中)。}。我这才第一次面对那场葬礼,不再是借文本的中介——那些文本固然用意味深长的言辞对我说着贬损、掠夺与对一部作品的屠杀,说着某种数学风格与某种数学进路(以缺席的大师为其人身)的埋葬——而是这一次,以一种直接而可触的方式,带着熟悉的面貌,借着一个熟悉的声音,音调殷勤而天真。葬礼终于就在我眼前,“有血有肉”,带着那副我如今已认得的殷勤而平淡的面貌;但这是我第一次用新的眼睛、新的注意去看它。于是,那个在我前几个月的反思中已显出是我隆重葬仪上的大主祭、是那位“披着祭袍的祭司”、同时又是一场史无前例的“手术”的主要匠人与主要“受益者”、一部被交付给取笑与掠夺的作品的隐形继承人的人,就这样在我面前展开……

这次会面正落在《收获与播种》里“第三道浪”的开端,当时我刚投身那场关于阴与阳的漫长冥想,追着一串飘忽而执着的联想。当时,这段短短的插曲只留下几行回声的痕迹,一笔带过。然而它标志着一个重要的时刻,它的果实要几个月后才清楚地显现。

还有第二次这样面对“有血有肉的葬礼”的时刻。那不过是十天前的事,它又一次“在最后一刻”重新启动了一场总也不断重新开张的调查。这一次,只是一个打给让–皮埃尔·塞尔(Jean-Pierre Serre)的电话\fnote{23}{}{这几乎是逐字引自按语《掘墓人——或整个教团》(第97号,第417页)。}。这场“漫无边际”的谈话,以触目的方式、甚至超出一切预料地证实了我(不过几天前)刚刚长篇大论地对自己讲明白的东西\fnote{24}{}{在同一则按语(第173号)的 c 部分(“那唯一者——或默许”)。},而且几乎是违我本意地讲明白的:关于塞尔在我的葬礼里所扮演的角色,关于他内心里对那件“就在他眼皮底下”发生的事的一种“秘密默许”,而他从不装作看见或感觉到什么。

这里同样,理所当然,谈话是再“酷”不过、再友好不过的;而且塞尔对我这种友好的态度显然也再真诚不过、再真实不过。这并不妨碍这一次我真正看见,或者——我真想写——“摸到”了我终于不得不承认的那种“默许”;它无疑是“秘密的”(如我先前所写),但尤其
% ---------------- p.89 ----------------
\origpage{L26}是殷勤的,正如我当时所能毫无疑义地看到的那样。一种殷勤而毫无保留的默许:让该被埋掉的被埋掉;并且,凡在看来合宜之处,无论用什么手段,都要让一种真实的(塞尔对此有第一手的了解)、不受欢迎的生父,被一个伪造而受欢迎的生父替换掉……\fnote{25}{}{这几乎是逐字引自按语《掘墓人——或整个教团》(第97号,第417页)。}。

\origpage{L27}\pmark 这触目地证实了一年前已经出现的一个直觉,当时我写道\fnote{26}{}{这段引文出自同一则按语(见上页脚注),同在第417页。}:

\begin{quote}
“从这一光照下看\fnote{27}{}{所谓“这一光照”,指的是刚才谈到的那种蓄意,即不惜一切代价除掉“不受欢迎的生父”(甚至“不可容忍的”,这里沿用所引按语中的说法)。},主祭德利涅便不再显得是那个按照决定他自身生命与行动的那些深刻力量、塑造出一种风尚的人,而更像是那“合法的继承人”\fnote{28}{}{德利涅这个“继承人”的角色,是一个既隐形(德利涅已发表的文字里,没有一行能让人猜到他可能从我的嘴里学到过什么),同时又清楚地被所有人感觉到、承认的角色。这正是德利涅两面手法及其“风格”的典型侧面之一:他精熟地在这含混之上做戏,收取这种心照不宣的继承人角色所带来的好处,同时又不认那位已故的师父,并主持着规模宏大的埋葬作业。}的指定工具,服务于一种具有毫无裂缝之一致性的集体意志——那意志执着于一项不可能的任务:把我和我个人的风格从当代数学里抹掉。”
\end{quote}

如果说德利涅当时在我眼里是那种“具有毫无裂缝之一致性的集体意志”的“指定工具”(同时也是头一个、主要的“受益者”),那么塞尔现在在我眼里,就是这同一种集体意志的化身,是它毫无保留之默许的担保人;那默许给予一切数不清的舞弊与欺诈,直至那些大规模的集体蒙骗与肆无忌惮的侵占“作业”,只要这些都有助于那项“不可能的任务”——针对我这点微不足道的、已故的人身,或者针对另一位\fnote{29}{}{我在这里想到的是左格曼·梅布胡特(Zoghman Mebkhout),第一次谈到他是在导论第6节(“葬礼”),然后是在按语《我的孤儿们》(第46号)里,以及在(较晚写成的、在发现葬礼之后的)按语《一场教学的失败(二)——或创造与自得》和《一种不义与无力的感觉》(第44$'$、44$''$号)里。我是在构成葬礼“第七支队列”的那十一则按语《讨论会——或梅布胡特的层与反常》(第75—80号)中,一点一点发现那桩对梅布胡特先驱性工作的不义抹除与侵占的。关于这(第四次、也是最后一次)“运算”,一份更详尽的调查与叙述,构成了“四种运算”这场调查里最充实的部分,用的正是那个理所当然的名字:《封神》(第171(i)至171号诸按语)。}竟敢自称出于我、并且不顾所有人反对而摆出“格罗滕迪克的继承者”架势的人。

\pmark 葬礼之中有许多悖谬而令人不知所措的侧面,其中之一便是:它首先是——甚至可以说是唯独是——由那些曾是我的朋友或学生的人做成的,而在那个世界里,我从不曾知道自己有敌人。我想,首先正是由于这一点,《收获与播种》才与你关系格外近,而我这封正在写的信也才要反过来成为一次招呼。因为如果你是数学家,如果你是他们中的一个——曾是我的学生,或曾是我的朋友——那么你与这场葬礼大概并非无关:或者以行动,或者以默契,哪怕只是以你对我的沉默——对一件就在你家门口台阶前进行的事保持沉默。而如果(说来稀奇)你接受我这些卑微的话、接受它们带给你的见证,而不是闭门不出、把送信人当不受欢迎的人打发走,那么你或许就会明白:被所有人、并且有你参与(积极的参与,或心照不宣的默许)埋掉的,不只是另一个人的作品——那是他与数学相爱的果实与活生生的见证;而且,在这个埋葬(它从不叫出自己的名字……)之下更深、更隐秘的层面上,被埋掉的是你自己存在中活着而本质的一部分,是你原初的认识、爱与创造的力量;你乐于亲手借着另一个人的身体把它埋掉。

% ---------------- p.90 ----------------
\origpage{L28}在我所有的学生里,德利涅占着一个很特殊的位置,关于这一点我在反思中谈得很多\fnote{30}{}{关于这一点,尤其见 RS II 里那十七则一组的按语《吾友皮埃尔》(第60—71号)。}。他曾经是——而且远远是——最“亲近”的一个,也是唯一一个(不论学生与否)把某种宏大的眼光亲密地吸收过来、并据为己有\fnote{31}{}{这种“宏大的眼光”,德利涅确确实实“亲密地吸收过来、并据为己有”了;它对他也曾有过强大的魅力,而且至今仍不由自主地魅惑着他,尽管同时有一股强横的力量推着他去摧毁它,去炸开它根本的统一,去把散落的碎片据为己有。于是,他对一位被否认的、“已故的”师父那种隐形的敌意,正是他自身内部分裂的表现;这分裂深深地印在他离去之后的作品上——那部作品远远停在我所了解的他那相当惊人的才能之下。}的人;那眼光在我身上诞生、长大,远在我俩相遇之前。而在我所有与我共有数学激情的朋友里,最亲近的是塞尔——他同时也多少算个长辈——(同样远远地)最亲近,尤其是因为他在十年间在我的工作里扮演了一个独一无二的角色:某些重大投入的“引爆器”,以及五十、六十年代启发我数学思考的多数有力观念(直到我离去之时)的“引爆器”。他们两人与我这个人的那种很特别的关系,当然与他们各自非凡的才能不无关联;这才能使他们对同代以及其后几代的数学家拥有同样非凡的威望。撇开这些共同点,塞尔与德利涅的性情与作风,在我看来却再相异不过,在许多方面彼此正相对立。

不管怎样,如果说有哪位数学家曾以某种名义“亲近”我这个人、我的作品(而且,更有甚者,是出了名地如此),那正是塞尔与德利涅:一位是长辈,是在某种眼光孕育的关键时期我作品中的灵感之源;另一位是我最有天分的学生,而他反过来把我当作(并且一直当作,不管有没有那场葬礼……)他主要的(也是秘密的……)灵感之源\fnote{32}{}{关于这一点,见上页脚注。}。如果在我离去之后不久,一场葬礼就开动起来(并且正儿八经地变成了“死亡”),并且落实为一支没完没了的、大大小小“作业”的队列,服务于同一个目的,那么这事只可能是在这两位——前长辈与前学生(甚至前“门徒”)——协同而紧密团结的配合下完成的:一位悄悄而有效地主持着作业,同时召集我某些学生归队\fnote{33}{}{这里说的,确确实实是另外那五个以簇的上同调为主题(与德利涅一样)的学生。},他们正巴不得屠杀父亲(以那个肥硕臃肿、嗡嗡饶舌的超级娘们的怪诞可笑形象为靶子);另一位则毫无保留、无条件、无限制地对(四项)作业的继续开“绿灯”(那些作业是:贬损、屠杀、宰割、分赃,分一具取之不尽的尸首……)。

\section{(9) 清空}

\origpage{L29}正如我刚才已经暗示过的,我必须克服相当大的内心抵抗,或者更确切地说,要靠一种耐心、细致而顽强的工作让它们自行消解,才能与某些熟悉的形象分开——那些形象牢牢盘踞、惰性极大,几十年来在我身上(正如在所有人身上,想必也在你身上)取代了对现实直接而细腻的感知,具体说来,是对某个数学世界的感知,而一段过去与一部作品仍然把我与那个世界连在一起。这些形象、这些现成观念里最根深蒂固的一条是:一个国际知名的学者,甚至一个看上去是大数学家的人,竟然会去干大大小小的欺诈,这从一开始就被排除在外(哪怕只是作为例外,更不必说作为一种亲爱的
% ---------------- p.91 ----------------
\origpage{L30}习惯……);或者,即便他(又出于老习惯)自己不伸手去沾,竟会张开双臂欢迎别人搞起来的那些“(有时简直不顾一切体面感)”的作业,而在其中,由于这样或那样的原因,他还能得利。

我身上这种精神的惰性竟到了这个地步:仅仅不到两个月前,在整整一年之久的漫长反思之后,我才终于怯生生地瞥见:塞尔在这葬礼里也许也有份——而这件事如今在我看来已是明证,甚至与最近那次意味深长的谈话无关。如同对那个曾在我起步时善意接纳我的“布尔巴基圈子”的所有成员一样,在他身上尤其如此,对我来说存在着某种心照不宣的“禁忌”,围绕着他这个人。他代表着某种“优雅”的化身——这种优雅绝不限于形式,它还包括一种严格、一种一丝不苟的正直。

在我去年4月19日发现那场葬礼之前,我做梦也不会想到:我某个旧日的学生竟会在执业中做出不诚实的事,无论是对我,还是对任何人;而如果这假设落在我最有天分的那一个、也最亲近的那一个身上,在我看来就最为荒诞不经!然而,早从我离去之时起,直到其后这些年、直到今天,我一直有充分机会看清他与我的关系是何等分裂。也不止一次,我见他(看来只为取乐)运用那种使人气馁并羞辱人的权力,只要时机合适。每一次我都深受触动(大概比我愿意向自己承认的更深……)。这些都是某种深刻失序的相当意味深长的迹象;而这失序(我有充分机会确认)绝不限于他一个人,即使在我旧日学生中最狭小的圈子里也是如此。这种失序,由于丧失对他人的敬重,其触目与深刻程度,并不亚于那种表现为所谓“职业上的不诚实”的失序。尽管如此,\dmark 发现这样一种不诚实,对我仍是一个十足的意外、一记重击。

在那场令人目瞪口呆的揭示之后——随后又是一整串同一路数的“瀑布”般的发现——的几周里,我渐渐明白:在我某些学生中间\fnote{34}{}{见上页脚注。},某种舞弊早在我离去之前的那些年里就已经开始了。恰恰在我最有天分的那一个身上,这一点格外触目——就是他,在我离去之后定了调子,并且(如我刚才所写)“悄悄而有效地主持着作业”。隔了将近二十年回看,这舞弊如今在我看来已是明摆着的事,简直“扎眼”。如果说当时我选择闭上眼不看正在发生的事,一心在那个“一切都只是秩序与美”的世界里(如我乐于想象的那样)追逐“白鲸”,那么我今天确认:我当时未能承担起我应负的责任——对那些在我身边学一门我所爱的行当的学生应负的责任;那行当还不只是一门手艺,也不只是养成某种“嗅觉”。由于对出色的学生徇情,由于我乐于(凭一道心照不宣的判决)把他们当作“特殊的存在”、当作无可怀疑的人,我当时就以我那一份\fnote{35}{}{那份“贡献”,尤其见于按语《特殊的存在》(第67$'$号),以及《攀升》与《含混》两则按语(第63$'$、63$''$号),后来又(在一种略有不同的光照下)见于按语《逐出》(第169号)末尾。另一类“贡献”见于《愚妄与更新》,那是对天分不那么出色的年轻数学家的自得态度。这种对普遍败坏中所负一份责任的认识,在《竞技数学》一节(第40号)里达到顶点。}促成了我今天看到一个世界、以及一些曾为我所爱的人中间蔓延开来的那种(在我看来史无前例的)腐败。

当然,鉴于那些东西惰性极大,要与之分离——与人们习惯称之为“幻觉”的东西分离(那语调里不无几分惋惜……),而我宁愿称之为
% 下接 part1-ch03b.tex
